\chapter{Kết luận}
\label{sec:ket_luan}

Việc số hóa hạ tầng viễn thông thông qua mô hình đám mây điểm 3D chính xác ngày càng quan trọng, đặc biệt trong quản lý thiết bị, bảo trì và vận hành trạm viễn thông. Tuy nhiên, bài toán tái tạo mây điểm 3D chi tiết cột ăng-ten gặp rất nhiều thách thức do khoảng cách chụp xa, bề mặt thiết bị trơn nhẵn và kết cấu cột lặp lại, ảnh hưởng đến ước tính camera pose và chất lượng của mô hình mây điểm 3D. Để giải quyết vấn đề này, luận văn đề xuất phương pháp cải tiến Hybrid-SfM, kết hợp thuật toán trích xuất, so khớp đặc trưng dựa trên học sâu ALIKED \cite{zhao2023aliked} và LightGlue \cite{lindenberger2023LightGlue} cho vùng ảnh chứa cột - được phát hiện bởi mô hình YOLO-v11 và thuật toán trích xuất đặc trưng truyền thống SIFT cho toàn bộ bức ảnh. Bên cạnh đó, việc thực hiện bước lọc dựa trên bản đồ độ sâu tương đối từ Depth Anything V2 giúp loại bỏ nhiễu, nâng cao độ ổn định, độ chính xác tính toán camera pose cho quy trình SfM tăng dần. Kết quả thực nghiệm trên 5 bộ dữ liệu ảnh cột thực tế cho thấy phương pháp Hybrid-SfM  đề xuất giúp giảm đáng kể sai số chiếu lại (MRE) so với phương pháp SfM truyền thống. Việc tạo ra đám mây điểm 3D chi tiết, chính xác cho cột ăng-ten trạm viễn thông là bước nền tảng quan trọng, tiến tới xây dựng một mô hình 3D trạm viễn thông hoàn chỉnh, trực quan và giàu thông tin. Mô hình này chính là tiền đề thiết yếu để có thể triển khai hiệu quả các hoạt động bảo dưỡng, bảo trì, kiểm tra định kỳ và giám sát chất lượng hạ tầng trạm trong thực tế. Cuối cùng, luận văn đưa ra những so sánh, đánh giá chất lượng tái tạo mô hình mây điểm 3D chi tiết bằng MVS (COLMAP MVS) và NeRF (biến thể nerfacto) cho thấy NeRF vượt trội trong phục dựng bề mặt trơn nhẵn, đồng màu, đạt mật độ điểm cao, phù hợp cho các tác vụ liên quan tới quản lí thiết bị vô tuyến trên cột sau này. Dù NeRF còn hạn chế về độ chính xác cấu trúc hình học, khả năng bao phủ toàn diện của nó là một ưu thế lớn, có tính khả thi khi ứng dụng vào bài toán thực tế này.

Mặc dù phương pháp đề xuất đã đạt được những kết quả khả quan, vẫn còn nhiều hướng tiềm năng để tiếp tục cải thiện những hạn chế và mở rộng nghiên cứu trong tương lai: (1) Nghiên cứu giải pháp tối ưu thời gian tính toán cho phương pháp Hybrid-SfM ; (2) Nghiên cứu cải tiến độ chính xác cấu trúc hình học cho mô hình 3D, bằng cách bổ sung các ràng buộc hình học cho mô hình NeRF trong quá trình huấn luyện, tái tạo. Từ đó, giúp tạo ra một mô hình 3D vừa có mật độ điểm, độ bao phủ tốt, vừa đảm bảo cấu trúc hình học chính xác cho những bài toán quản lý hạ tầng và thiết bị vô tuyến. Những hướng nghiên cứu này hứa hẹn sẽ tiếp tục nâng cao tính khả thi ứng dụng giải pháp vào các bài toán thực tế cho không chỉ lĩnh vực viễn thông mà còn rất nhiều các lĩnh vực liên quan khác trong tương lai.